\documentclass[hyperref={pdfpagelabels=false}]{beamer}

\usepackage{ragged2e}
\renewcommand{\raggedright}{\leftskip=0pt \rightskip=0pt plus 0cm}

\usepackage{listings}
\usepackage{lmodern}
\usepackage{fontspec}
\usepackage{xeCJK}
\setCJKmainfont{Noto Sans CJK TC}
\setCJKsansfont{Noto Sans CJK TC}
\setCJKmonofont{Noto Sans Mono CJK TC}

\usepackage[english]{babel}
\usepackage{amsmath,amssymb,amsfonts,amsthm}

\hypersetup{
  colorlinks=true,
  urlcolor=blue,
  citecolor=blue,
  pdfborder={0 0 0}
}

\usepackage{accsupp}
\newcommand{\emptyaccsupp}[1]{\BeginAccSupp{ActualText={}}#1\EndAccSupp{}}

\usepackage{tikz}
\usetikzlibrary{arrows,decorations.pathmorphing,backgrounds,fit,positioning,shapes.symbols,chains}

\usepackage{booktabs}
\usepackage{threeparttable}
\usepackage{colortbl}
\usepackage{xcolor}
\usetikzlibrary{tikzmark}

\newcommand{\indep}{\perp\!\!\!\perp}
\newcommand{\nindep}{\not\!\perp\!\!\!\perp}

\beamertemplatefootpagenumber

\lstdefinelanguage{Stata}{
	morekeywords=,
	morekeywords=[2]{cntrade, chinafin, wbopendata, spmap,},
	morekeywords=[3]{regress, summarize, display,},
	morekeywords=[4]{forvalues, if, foreach, set},
	morekeywords=[5]{rnormal, runiform},
	morecomment=[l]{//},
	morecomment=[s]{/*}{*/},
	morecomment=[n][keywordstyle9]{`}{'},
	morestring=[b]",
	sensitive=true,
}

\lstdefinestyle{numbers}{
	numbers=left,
	stepnumber=1,
	numberstyle=\tiny\emptyaccsupp,
	xleftmargin=2em,
}

\lstdefinestyle{nonumbers}{
	numbers=none,
}

\lstdefinestyle{stata-plain}{
	language=Stata,
	basicstyle=\footnotesize\ttfamily,
}

\lstdefinestyle{stata-editor}{
	language=Stata,
	basicstyle=\footnotesize\ttfamily,
	keywordstyle={\color{NavyBlue}},
	keywordstyle=[2]{\color{NavyBlue}},
	keywordstyle=[3]{\color{NavyBlue}},
	keywordstyle=[4]{\color{NavyBlue}},
	keywordstyle=[5]{\color{Blue}},
	keywordstyle=[9]{\color{TealBlue}},
	stringstyle=\color{Maroon},
	commentstyle=\color{OliveGreen},
}

\lstset{
	language=Stata,
	style=stata-plain,
	style=numbers,
	showstringspaces=false,
	breaklines,
	frame=single,
}


\title[]
{AI Tools for Empirical Research}

\author
{Prof. Tzu-Ting Yang \\ 楊子霆}

\institute{Institute of Economics, Academia Sinica \\ 中央研究院經濟研究所}

\date{\today}


\usetheme{CambridgeUS}
\setbeamertemplate{footline}{%
	\hfill\usebeamercolor[fg]{page number in head/foot}%
	\usebeamerfont{page number in head/foot}%
	\insertframenumber\,/\,\inserttotalframenumber\hspace{1em}\vspace{2pt}%
}

\beamersetuncovermixins{\opaqueness<1>{25}}{\opaqueness<2->{15}}

\begin{document}

% ------------------------------------------------------------------
%  TITLE PAGE
% ------------------------------------------------------------------

\begin{frame}{}
\titlepage
\end{frame}


% ------------------------------------------------------------------
%  SLIDE 1: Why Use AI Tools?
% ------------------------------------------------------------------

\begin{frame}{Why Use AI Tools in Empirical Research?}
\begin{columns}[T]

\begin{column}{0.47\textwidth}
\textbf{\textcolor{red}{Traditional Workflow}}
\vspace{0.2cm}
\begin{itemize}
	\item Manually read long codebook PDFs
	\vspace{0.15cm}
	\item Guess variable meanings from short names
	\vspace{0.15cm}
	\item Write data cleaning code from scratch
	\vspace{0.15cm}
	\item Debug cryptic error messages alone
	\vspace{0.3cm}
	\item[\textcolor{red}{$\Rightarrow$}] \textbf{\textcolor{red}{Time-consuming \& error-prone}}
\end{itemize}
\end{column}

\begin{column}{0.47\textwidth}
\textbf{\textcolor{teal}{With Claude Code}}
\vspace{0.2cm}
\begin{itemize}
	\item Ask AI to summarize key variables
	\vspace{0.15cm}
	\item Generate variable labels \& descriptions
	\vspace{0.15cm}
	\item Auto-generate cleaning \& merging code
	\vspace{0.15cm}
	\item Debug errors with AI explanation
	\vspace{0.3cm}
	\item[\textcolor{teal}{$\Rightarrow$}] \textbf{\textcolor{teal}{Focus on research design}}
\end{itemize}
\end{column}

\end{columns}

\end{frame}


% ------------------------------------------------------------------
%  SLIDE 2: What is Claude Code?
% ------------------------------------------------------------------

\begin{frame}{What is Claude Code?}

\begin{itemize}
	\item \textbf{Claude Code} is Anthropic's AI coding assistant
	\vspace{0.2cm}
	\begin{itemize}
		\item Reads and writes your actual project files (do-files, R scripts, datasets)
		\vspace{0.15cm}
		\item Understands your codebook, error messages, and Stata/R output
		\vspace{0.15cm}
		\item Generates, edits, and debugs code through a conversation
	\end{itemize}
	\vspace{0.3cm}
	\item \textbf{What it can do for empirical research:}
	\vspace{0.15cm}
	\begin{itemize}
		\item Explain what each variable means from a codebook
		\vspace{0.15cm}
		\item Write Stata or R code for data cleaning, sample construction, and analysis
		\vspace{0.15cm}
		\item Translate a Stata do-file into R (or vice versa)
		\vspace{0.15cm}
		\item Explain what a command does and why an error occurred
		\vspace{0.15cm}
		\item Suggest how to construct variables for a specific research design
	\end{itemize}
\end{itemize}

\end{frame}


% ------------------------------------------------------------------
%  SLIDE 3: What You Need to Install
% ------------------------------------------------------------------

\begin{frame}{Getting Started: What to Install First}

\begin{itemize}
	\item The following tools are needed for \textbf{both} modes of Claude Code:
	\vspace{0.3cm}
	\item \textbf{Claude Desktop App}
	\vspace{0.1cm}
	\begin{itemize}
		\item \url{https://claude.ai/download}
		\vspace{0.1cm}
		\item Requires an Anthropic account (Pro plan to use Claude Code mode)
	\end{itemize}
	\vspace{0.3cm}
	\item \textbf{Git} --- version control; Claude Code uses it to track file changes
	\vspace{0.1cm}
	\begin{itemize}
		\item \url{https://git-scm.com/downloads}
	\end{itemize}
	\vspace{0.3cm}
	\item \textbf{VS Code} --- recommended code editor
	\vspace{0.1cm}
	\begin{itemize}
		\item \url{https://code.visualstudio.com}
		\vspace{0.1cm}
		\item Useful for viewing and editing your do-files and R scripts alongside Claude Code
	\end{itemize}
\end{itemize}

\end{frame}


% ------------------------------------------------------------------
%  SLIDE 4: Two Ways to Use Claude Code
% ------------------------------------------------------------------

\begin{frame}[fragile]{Two Ways to Use Claude Code}

\begin{itemize}
	\item \textbf{Option 1: Desktop App} \hfill {\small (recommended for beginners)}
	\vspace{0.1cm}
	\begin{itemize}
		\item Open the Claude desktop app and switch to \textbf{Claude Code mode}
		\vspace{0.1cm}
		\item Graphical interface --- easier to get started
		\vspace{0.1cm}
		\item Can read and write your project files, just like the terminal version
	\end{itemize}
	\vspace{0.35cm}
	\item \textbf{Option 2: Terminal} \hfill {\small (more flexible for advanced users)}
	\vspace{0.1cm}
	\begin{itemize}
		\item Run from the command line directly inside your project folder
		\vspace{0.1cm}
		\item On Windows: open \textbf{PowerShell}; on Mac/Linux: open \textbf{Terminal}
		\vspace{0.1cm}
		\item First install Node.js if not already present: \url{https://nodejs.org}
		\vspace{0.1cm}
		\item Then install Claude Code by typing:
\begin{lstlisting}[style=nonumbers, basicstyle=\small\ttfamily, frame=single]
npm install -g @anthropic-ai/claude-code
\end{lstlisting}
		\item Navigate to your project folder and launch with:
\begin{lstlisting}[style=nonumbers, basicstyle=\small\ttfamily, frame=single]
claude
\end{lstlisting}
	\end{itemize}
\end{itemize}

\end{frame}


% ------------------------------------------------------------------
%  SLIDE 5: Free Alternative --- Gemini CLI
% ------------------------------------------------------------------

\begin{frame}[fragile]{Free Alternative: Google Gemini CLI}

\begin{itemize}
	\item \textbf{Gemini CLI} is Google's open-source AI coding agent for the terminal
	\vspace{0.15cm}
	\begin{itemize}
		\item Reads and writes your project files, just like Claude Code
		\vspace{0.1cm}
		\item Powered by Gemini 3 Pro with a \textbf{1 million token} context window
		\vspace{0.1cm}
		\item \textbf{Free tier}: 60 requests/min, 1,000 requests/day --- just use a Google account
	\end{itemize}
	\vspace{0.3cm}
	\item \textbf{Installation} (Windows: PowerShell; Mac/Linux: Terminal)
	\vspace{0.1cm}
	\begin{itemize}
		\item Step 1: Install Node.js (version 20 or higher): \url{https://nodejs.org}
		\vspace{0.1cm}
		\item Step 2: Install Gemini CLI globally:
\begin{lstlisting}[style=nonumbers, basicstyle=\small\ttfamily, frame=single]
npm install -g @google/gemini-cli
\end{lstlisting}
		\item Step 3: Navigate to your project folder and launch:
\begin{lstlisting}[style=nonumbers, basicstyle=\small\ttfamily, frame=single]
cd C:\your\project\folder
gemini
\end{lstlisting}
		\item Step 4: First launch asks you to \textbf{Login with Google} --- select your account in the browser
	\end{itemize}
	\vspace{0.2cm}
	\item After login, type your task in plain language and Gemini reads your files and edits them
\end{itemize}

\end{frame}


% ------------------------------------------------------------------
%  SLIDE 6: Gemini Code Assist in VS Code
% ------------------------------------------------------------------

\begin{frame}{Using Gemini in VS Code: Gemini Code Assist}

\begin{itemize}
	\item \textbf{Gemini Code Assist} is the VS Code extension that embeds Gemini in your editor
	\vspace{0.15cm}
	\begin{itemize}
		\item Chat with Gemini directly in the VS Code sidebar
		\vspace{0.1cm}
		\item \textbf{Agent mode}: give a task, Gemini proposes a plan and edits multiple files
		\vspace{0.1cm}
		\item Also free with a personal Google account
	\end{itemize}
	\vspace{0.3cm}
	\item \textbf{Installation steps}
	\vspace{0.1cm}
	\begin{itemize}
		\item Step 1: Open VS Code and go to the \textbf{Extensions} panel (Ctrl+Shift+X)
		\vspace{0.1cm}
		\item Step 2: Search for \textbf{``Gemini Code Assist''} and click Install
		\vspace{0.1cm}
		\item Step 3: Click the Gemini icon in the sidebar and \textbf{Sign in with Google}
		\vspace{0.1cm}
		\item Step 4: Open your project folder in VS Code, then start chatting or switch to \textbf{Agent mode}
	\end{itemize}
	\vspace{0.3cm}
	\item \textbf{Gemini CLI} and \textbf{Gemini Code Assist} share the same free quota ---\\
	      both run on the same Google account; no extra sign-up needed
\end{itemize}

\end{frame}


% ------------------------------------------------------------------
%  SLIDE 7: Setting Up VS Code for Stata
% ------------------------------------------------------------------

\begin{frame}[fragile]{Setting Up VS Code for Stata}

\begin{itemize}
	\item \textbf{Step 1: Install two extensions} (Ctrl+Shift+X to open Extensions panel)
	\vspace{0.1cm}
	\begin{itemize}
		\item \textbf{Stata Enhanced} (by Kyle Barron) --- syntax highlighting for .do files
		\vspace{0.1cm}
		\item \textbf{stataRun} (by Yeaoh) --- sends code to your running Stata application
	\end{itemize}
	\vspace{0.25cm}
	\item \textbf{Step 2: Configure stataRun} --- click the gear icon next to the extension
	\vspace{0.1cm}
	\begin{itemize}
		\item Set \texttt{stataRun.whichApp} to match your version: \texttt{stataMP}, \texttt{stataSE}, etc.
		\vspace{0.1cm}
		\item On Windows, also verify the Stata installation path is correct
	\end{itemize}
	\vspace{0.25cm}
	\item \textbf{Step 3: Open your project folder} in VS Code (File $\to$ Open Folder)
	\vspace{0.1cm}
	\begin{itemize}
		\item Open a .do file --- it will now display with colour syntax
		\vspace{0.1cm}
		\item Make sure Stata is already running in the background
	\end{itemize}
	\vspace{0.25cm}
	\item \textbf{Key shortcuts} (stataRun defaults):
	\vspace{0.1cm}
	\begin{itemize}
		\item Ctrl+Shift+S \hspace{0.3cm} Run selected lines (or current line)
		\vspace{0.05cm}
		\item Ctrl+Shift+A \hspace{0.3cm} Run entire do-file
		\vspace{0.05cm}
		\item Ctrl+Shift+D \hspace{0.3cm} Run from current line to end
	\end{itemize}
\end{itemize}

\end{frame}


% ------------------------------------------------------------------
%  SLIDE 8: Setting Up VS Code for R
% ------------------------------------------------------------------

\begin{frame}[fragile]{Setting Up VS Code for R}

\begin{itemize}
	\item \textbf{Step 1: Install R} from \url{https://cran.r-project.org} if not already installed
	\vspace{0.2cm}
	\item \textbf{Step 2: Install two R packages} (run in R or RStudio):
\begin{lstlisting}[style=nonumbers, basicstyle=\small\ttfamily, frame=single]
install.packages("languageserver")  # autocomplete & diagnostics
install.packages("httpgd")          # plot viewer inside VS Code
\end{lstlisting}
	\vspace{0.15cm}
	\item \textbf{Step 3: Install the R extension in VS Code}
	\vspace{0.1cm}
	\begin{itemize}
		\item Ctrl+Shift+X $\to$ search \textbf{``R''} (by REditorSupport) $\to$ Install
		\vspace{0.1cm}
		\item If VS Code cannot find R automatically: go to extension settings $\to$ \texttt{R: Rpath} $\to$ paste the path to your R executable
	\end{itemize}
	\vspace{0.2cm}
	\item \textbf{Step 4: Open an .R file and start coding}
	\vspace{0.1cm}
	\begin{itemize}
		\item Ctrl+Enter \hfill Run current line or selection
		\vspace{0.05cm}
		\item Ctrl+Shift+S \hfill Source (run) the entire file
		\vspace{0.05cm}
		\item Plots appear in the \textbf{Viewer} panel on the right (powered by httpgd)
	\end{itemize}
\end{itemize}

\end{frame}


% ------------------------------------------------------------------
%  SLIDE 9: Combining VS Code with Gemini and Claude
% ------------------------------------------------------------------

\begin{frame}{Combining VS Code with Gemini and Claude}

\begin{itemize}
	\item Once Stata and R are working in VS Code, add an AI assistant in the same window
	\vspace{0.3cm}
	\item \textbf{Option A: Gemini Code Assist} \hfill {\small (free)}
	\vspace{0.1cm}
	\begin{itemize}
		\item Extensions $\to$ search \textbf{``Gemini Code Assist''} $\to$ Install
		\vspace{0.1cm}
		\item Sign in with Google account in the sidebar $\to$ ready to use
		\vspace{0.1cm}
		\item In the chat panel: paste your error, ask Gemini to fix your .do or .R file
		\vspace{0.1cm}
		\item Switch to \textbf{Agent mode} to let Gemini read \& edit files autonomously
	\end{itemize}
	\vspace{0.3cm}
	\item \textbf{Option B: Claude Code} (terminal mode) \hfill {\small (Pro plan required)}
	\vspace{0.1cm}
	\begin{itemize}
		\item Open the \textbf{integrated terminal} in VS Code: Ctrl+` (backtick)
		\vspace{0.1cm}
		\item Navigate to your project folder and type \texttt{claude}
		\vspace{0.1cm}
		\item Claude reads all files in the folder --- ask it to write, fix, or explain your code
		\vspace{0.1cm}
		\item Edit the file in VS Code, run it with stataRun/R, all in the same window
	\end{itemize}
	\vspace{0.2cm}
	\item Both tools can see your \textbf{.do files, .R scripts, and output logs}\\
	      $\Rightarrow$ describe your research question and let AI draft the code for you
\end{itemize}

\end{frame}


\end{document}